% ============================================================
% Problema 9
% ============================================================
\section{Problema 9: Capa 3: RIP no está aprendiendo rutas}

\begin{tcolorbox}[enunciado]
Se tienen tres routers: R1, R2 y R3. Los routers utilizan RIPv2 para intercambiar rutas.

\vspace{10pt}
\begin{center}
\begin{tikzpicture}[
    node distance=3.5cm, 
    texto/.style={align=center}
]
    % Definición de los Routers
    \node (R1) {$R1$};
    \node (R2) [right=of R1] {$R2$};
    \node (R3) [right=of R2] {$R3$};
    
    % Definición de las redes LAN 
    \node (L1) [above=0.4cm of R1, texto] {LAN A \\ 192.168.1.0/24};
    \node (L2) [above=0.4cm of R2, texto] {LAN B \\ 192.168.2.0/24};
    \node (L3) [above=0.4cm of R3, texto] {LAN C \\ 192.168.3.0/24};
    
    % Líneas de conexión (Cables)
    \draw (R1) -- (R2);
    \draw (R2) -- (R3);
    \draw (R1) -- (L1);
    \draw (R2) -- (L2);
    \draw (R3) -- (L3);
\end{tikzpicture}
\end{center}
\vspace{10pt}

La configuración esperada de RIP es:

\vspace{10pt}
\begin{center}
\begin{minipage}{5.5cm}
\begin{verbatim}
R1: network 192.168.1.0
    network 10.0.12.0

R2: network 10.0.12.0
    network 10.0.23.0
    network 192.168.2.0

R3: network 10.0.23.0
    network 192.168.3.0
\end{verbatim}
\end{minipage}
\end{center}
\vspace{10pt}

Un equipo de la LAN A puede comunicarse con R1, pero no puede llegar a un equipo de la LAN C.

Al ejecutar \texttt{show ip route} en R1, no aparece ninguna ruta hacia 192.168.3.0/24.

\begin{enumerate}[label=\arabic*., leftmargin=*]
    \item ¿A qué capa del modelo OSI corresponde principalmente este problema?
    \item ¿Qué protocolo de enrutamiento se está utilizando?
    \item ¿Qué comando utilizarías para comprobar si RIP está funcionando?
    \item ¿Qué comando utilizarías para revisar las rutas aprendidas?
    \item ¿Qué podría ocurrir si R3 tiene mal configurado el comando network?
    \item ¿Cómo verificarías la conectividad entre los routers?
    \item ¿Qué ruta debería aparecer en R1 para llegar a 192.168.3.0/24?
\end{enumerate}
\end{tcolorbox}

\subsection*{Solución}
% Escribe aquí tu solución para el Problema 9...